%%%%%%%%%%%%%%%%%%%%%%%%%%%%%%%%%%%%%%%%%%%%%%%%%%%%%%%%%%%%%%%%%%%%%%
% Lớp 11 | Hàm số lượng giác và phương trình lượng giác | Phương trình lượng giác cơ bản
% Dạng 1D1?5-2: Điều kiện có nghiệm
% Số câu: 47  (sinh tự động bởi tach_cau.py; từ nay thầy sửa trực tiếp file này)
%%%%%%%%%%%%%%%%%%%%%%%%%%%%%%%%%%%%%%%%%%%%%%%%%%%%%%%%%%%%%%%%%%%%%%

%%%%%%%%%%%%%%%% TRẮC NGHIỆM 4 PHƯƠNG ÁN %%%%%%%%%%%%%%%%

% nguồn: Ok_Toan11_TracNghiem_Lan2
\begin{ex}%[1D1N5-2]%[ma:f2850f62]
	Có bao  nhiêu giá trị nguyên của tham số $m$ để phương trình $2\sin 2x+7-m=0$ có nghiệm?
	 \choice
		{$4$} 
		{\True  $5$}
		 {$6$}
		 {Vô số}
	\loigiai{ 
		Ta có: $2\sin 2x+7-m=0 \Leftrightarrow \sin 2x=\dfrac{m-7}{2}$.\\ 
		 Do đó phương trình có nghiệm $\Leftrightarrow-1\le \dfrac{m-7}{2} \le 1\Leftrightarrow 5\le m\le 9$ $\Rightarrow m\in \left\{5;6;7;8;9\right\}$.
	 }
\end{ex}

% nguồn: Ok_Toan11_TracNghiem_Lan2
\begin{ex}%[1D1N5-2]%[ma:81c0de23]
	Gọi $S$ là tập hợp  tất cả các giá trị nguyên của tham số $m$ để phương trình\\
	$\cos \left(2x-\dfrac{\pi}{3} \right)-m=2$ có nghiệm. Tính tổng $T$ của các phần tử  trong $S$.
	\choice 
		{$T=6$} 
		{$T=3$} 
		{$T=-2$} 
		{\True  $T=-6$}
	\loigiai{ 
		Phương  trình $\cos \left(2x-\dfrac{\pi}{3} \right)-m=2\Leftrightarrow \cos \left(2x-\dfrac{\pi}{3} \right)=m+2$.\\
		 Phương trình có nghiệm $\Leftrightarrow-1\le m+2\le 1\Leftrightarrow-3\le m\le-1$.\\
		Do $m\in \mathbb{Z}$ nên $S=\left\{-3;-2;-1\right\}$.\\ 
		 Vậy $T=\left(-3\right)+\left(-2\right)+\left(-1\right)=-6$.
	 }
\end{ex}

% nguồn: Ok_Toan11_TracNghiem_Lan2
\begin{ex}%[1D1N5-2]%[ma:51828b83]
	Phương  trình $m\cos x-1=0$ có nghiệm khi $m$ thỏa mãn điều kiện
	\choice 
		{$\heva{&m\ge-1  \\&m\le 1}$}
		 {\True $\hoac{&m\le-1\\&m\ge 1}$}
		 {$m\ge-1$}
		{$m\ge 1$} 
	\loigiai{ 
		Dễ thấy với $m=0$ thì phương  trình đã cho vô nghiệm.\\
		Với $m\ne 0$, ta có: $m\cos x-1=0\Leftrightarrow  \cos x=\dfrac{1}{m}$.\\
		 Phương trình đã cho có nghiệm $\Leftrightarrow$ $\left|\dfrac{1}{m} \right|\le 1\Leftrightarrow \dfrac{1}{\left|m\right|} \le 1\Leftrightarrow \left|m\right|\ge 1\Leftrightarrow \hoac{&m\ge 1 \\ &m\le-1}$.
	 }
\end{ex}

% nguồn: Ok_Toan11_TracNghiem_Lan2
\begin{ex}%[1D1N5-2]%[ma:5d403a30]
	Tìm tất cả các giá trị thực của tham số $m$ để phương  trình $\sin x=m$ có nghiệm thực.
	 \choice
		 {$m\ge 0$}
		{$m  < 0$}
		{$-1< m <  1$}
		{\True $-1\le m\le 1$} 
	 \loigiai{
		Do $-1\le \sin x\le 1, \; \forall  x\in \mathbb{R}$ nên phương trình $\sin x=m$ có nghiệm khi và chỉ khi $-1\le m\le 1$.
	 }
\end{ex}

% nguồn: Ok_Toan11_TracNghiem_Lan2
\begin{ex}%[1D1N5-2]%[ma:77d77641]
	Tìm giá trị thực của tham số  $m$ để phương trình $\left(m-2\right)\sin 2x=m+1$ nhận $x=\dfrac{\pi}{12}$ làm nghiệm.
	 \choice
		 {$m\ne 2$}
		{$m=\dfrac{2\left(\sqrt{3}+1\right)}{\sqrt{3}-2}$} 
		 {\True $m=-4$}
		{$m=-1$} 
	 \loigiai{
		Vì $x=\dfrac{\pi}{12}$ là một nghiệm của phương trình $\left(m-2\right)\sin  2x=m+1$ nên ta có\\
		 $\left(m-2\right)\sin \dfrac{2\pi}{12}=m+1\Leftrightarrow \dfrac{m-2}{2}=m+1\Leftrightarrow m-2=2m+2\Leftrightarrow m=-4$.\\
		Vậy $m=-4$ là  giá trị cần tìm.
	} 
\end{ex}

% nguồn: Ok_Toan11_TracNghiem_Lan2
\begin{ex}%[1D1N5-2]%[ma:043a7b08]
	Cho phương trình $\cos 5x=3m-5$. Gọi đoạn $\left[a;b\right]$ là tập hợp tất cả các giá trị của $m$  để phương trình có nghiệm. Tính $3a+b$.
	 \choice
		 {$5$}
		 {$-2$}
		{$\dfrac{19}{3}$} 
		{\True  $6$}
	 \loigiai{
		Phương trình đã cho có nghiệm khi  $-1\le 3m-5\le 1\Leftrightarrow 4\le 3m\le 6\Leftrightarrow \dfrac{4}{3} \le m\le 2$.\\
		Khi đó tập hợp tất  cả các giá trị của $m$ để phương trình có nghiệm là $\left[\dfrac{4}{3};2\right]$.\\
		Ta thu được  $a=\dfrac{4}{3}$; $b=2$.\\
		Suy  ra $3a+b=6$.
	} 
\end{ex}

% nguồn: Ok_Toan11_TracNghiem_Lan2
\begin{ex}%[1D1N5-2]%[ma:c614374a]
	Tìm m để phương trình $\sqrt{2} \sin \left(x+\dfrac{\pi}{4} \right)=m$  có nghiệm $x\in \left(0;\dfrac{\pi}{2} \right)$.
	 \choice
		 {$\hoac{&m > 1 \\&m\le \sqrt{2}}$}
		{$1\le m\le  \sqrt{2}$}
		{$1\le m <  \sqrt{2}$}
		{\True $1< m\le \sqrt{2}$} 
	\loigiai{ 
		Vì $0< x  < \dfrac{\pi}{2} \Rightarrow \dfrac{\pi}{4} < x+\dfrac{\pi}{4} < \dfrac{3\pi}{4} \Rightarrow \dfrac{\sqrt{2}}{2} < \sin \left(x+\dfrac{\pi}{4} \right)\le 1$.\\
		Phương trình đã cho có nghiệm $x\in \left(0;\dfrac{\pi}{2} \right)$ khi $\dfrac{\sqrt{2}}{2} < \dfrac{m}{\sqrt{2}} \le 1$  $\Leftrightarrow 1< m\le \sqrt{2}$.
	} 
\end{ex}

% nguồn: Ok_Toan11_TracNghiem_Lan2
\begin{ex}%[1D1N5-2]%[ma:9e86ca29]
	Có bao  nhiêu giá trị nguyên của tham số $m$ để phương trình $\cos x=m+1$ có nghiệm?
	 \choice
		{Vô số} 
		 {$1$}
		{$2$} 
		{\True $3$} 
	\loigiai{ 
		Phương trình $\cos x=m+1$ có nghiệm $\Leftrightarrow-1\le m+1\le  1\; \Leftrightarrow-2\le m\le 0$.\\
		Mà $m\in \mathbb{Z}\Rightarrow m\in \left\{-2;-1;0\right\}$.\\ 
		Vậy có $3$ giá  trị nguyên của tham số $m$ thoả mãn yêu cầu bài toán.
	 }
\end{ex}

% nguồn: Ok_Toan11_TracNghiem_Lan2
\begin{ex}%[1D1N5-2]%[ma:8505bc36]
	Có bao nhiêu số nguyên  $m$ để phương trình $3\cos \left(x+\dfrac{\pi}{6} \right)-m+5=0$ có nghiệm?
	 \choice
		 {$5$}
		{\True  $7$}
		 {$6$}
		 {$9$}
	\loigiai{ 
		Ta có: $3\cos \left(x+\dfrac{\pi}{6} \right)-m+5=0\Leftrightarrow \cos  \left(x+\dfrac{\pi}{6} \right)=\dfrac{m-5}{3}$.\\
		Điều kiện để phương trình có nghiệm: $-1\le \dfrac{m-5}{3} \le  1\Leftrightarrow 2\le m\le 8$.\\
		Do $m$ nguyên nên $m \in  \left\{2;3;4;5;6;7;8\right\}$.\\
		Vậy có 7 số nguyên $m$ thoả mãn yêu cầu bài toán. 
	} 
\end{ex}

% nguồn: Ok_Toan11_TracNghiem_Lan2
\begin{ex}%[1D1N5-2]%[ma:a14c5502]
	Tìm tất cả các giá trị của tham số $m$  để phương trình $\left(m+1\right)\sin x+2-m=0$ có nghiệm.
	 \choice
		{$m\le-1$} 
		{\True  $m\ge \dfrac{1}{2}$}
		{$-1< m\le  \dfrac{1}{2}$}
		{$m >-1$} 
	 \loigiai{
		Phương trình  $\left(m+1\right)\sin x+2-m=0\Leftrightarrow \left(m+1\right)\sin x=m-2\Leftrightarrow \sin x=\dfrac{m-2}{m+1}$.\\
		Để phương trình có  nghiệm
		\begin{align*}
			&-1\le \dfrac{m-2}{m+1}  \le 1\\
			&\Leftrightarrow  \heva{&0\le 1+\dfrac{m-2}{m+1} \\ &\dfrac{m-2}{m+1}-1\le 0}\\
			&\Leftrightarrow \heva{&\dfrac{2m-1}{m+1} \ge 0 \\&-\dfrac{3}{m+1} \le  0}\\
			&\Leftrightarrow  \heva {&\hoac {&m\ge \dfrac{1}{2} \\ & m < -1} \\ &m > -1}\\
			 &\Leftrightarrow m\ge \dfrac{1}{2}.
		\end{align*}
		Vậy $m\ge  \dfrac{1}{2}$ là giá trị cần tìm.
	} 
\end{ex}

% nguồn: Ok_Toan11_TracNghiem_Lan2
\begin{ex}%[1D1N5-2]%[ma:7980a5b8]
	Tìm tất cả các giá trị thực của tham số $m$  để phương trình $\cos x-m=0$ vô nghiệm.
	 \choice
		 {\True $m\in \left(-\infty;-1\right)\cup \left(1;+\infty \right)$}
		{$m\in (-\infty;-1]\cup [1;+\infty)$} 
		{$m\in \left(1;+\infty \right)$} 
		{$m\in (-\infty;-1)$} 
	 \loigiai{
		Do $\left|\cos x\right|\le 1$, $\forall  x\in \mathbb{R}$ nên phương trình: $\cos x-m=0\Leftrightarrow \cos x=m$
		 có nghiệm khi $\left|m\right|\le 1$ và vô nghiệm khi $\left|m\right|> 1$.
	} 
\end{ex}

% nguồn: Ok_Toan11_TracNghiem_Lan2
\begin{ex}%[1D1N5-2]%[ma:b49ff78c]
	Cho phương trình $\cos \left(2x-\dfrac{\pi}{3} \right)-m=2$. Tìm  $m$ để phương trình có nghiệm.
	\choice 
		{Không tồn  tại $m$}
		 {$m\in \left[-1;3\right]$}
		{\True $m\in  \left[-3;-1\right]$}
		{$m\in \mathbb{R}$} 
	\loigiai{ 
		Ta có $\cos  \left(2x-\dfrac{\pi}{3} \right)-m=2$ $\Leftrightarrow$ $\cos \left(2x-\dfrac{\pi}{3} \right)=m+2$.\\
		Lại có $-1\le  \cos \left(2x-\dfrac{\pi}{3} \right)\le 1$, nên phương trình có nghiệm khi
		\begin{align*}
			&-1\le m+2\le 1\\  &\Leftrightarrow-3\le m\le-1.
		\end{align*}
	 }
\end{ex}

% nguồn: Ok_Toan11_TracNghiem_Lan2
\begin{ex}%[1D1N5-2]%[ma:d8e9f5aa]
	Tìm  $m$ để phương trình $\cos x-2m+1=0$ có nghiệm.
	 \choice
		 {$m >-\dfrac{1}{2}$}
		{$0< m <  1$}
		{\True  $0\le m\le 1$}
		{$m\ge-\dfrac{1}{2}$} 
	 \loigiai{
		Ta có $\cos x-2m+1=0$ $\Leftrightarrow \cos x=2m-1$ có nghiệm khi và  chỉ khi
		\begin{align*}
			&-1\le 2m-1\le 1\\ 
			&\Leftrightarrow  0\le 2m\le 2\\
			&\Leftrightarrow  0\le m\le 1.
		\end{align*}
	} 
\end{ex}

% nguồn: Ok_Toan11_TracNghiem_Lan2
\begin{ex}%[1D1N5-2]%[ma:dc4380cd]
	Có bao nhiêu giá trị nguyên của tham  số $m$ để phương trình $3\sin 2x-m^2+5=0$ có nghiệm?
	 \choice
		 {$6$}
		{\True  $2$}
		{$1$} 
		{$7$} 
	\loigiai{ 
		Phương trình đã cho tương  đương với phương trình $\sin 2x=\dfrac{m^2-5}{3}$.\\
		 Vì $\sin 2x\in \left[-1;1\right]$ nên $\dfrac{m^2-5}{3} \in \left[-1;1\right]\Leftrightarrow m^2\in \left[2;8\right]\Leftrightarrow \hoac {&{-2\sqrt{2} \le m\le-\sqrt{2}} \\&{\sqrt{2} \le m\le 2\sqrt{2}}
		 .}$\\
		 Vậy có $2$ giá trị nguyên của của tham số $m$ thoả mãn yêu cầu bài toán.
	} 
\end{ex}

% nguồn: Ok_Toan11_TracNghiem_Lan2
\begin{ex}%[1D1N5-2]%[ma:45a04c6d]
	Tìm $m$ để phương trình $\cos 2x=m-1$ có nghiệm. 
	\choice 
		{\True $0\le m\le  2$}
		{$-1\le  m\le 1$}
		{$-2\le m\le 2$} 
		{$0\le m\le 1$} 
	\loigiai{ 
		Phương trình $\cos 2x=m-1$ có nghiệm  khi và chỉ khi $-1\le m-1\le 1 \Leftrightarrow 0\le m \le 2$.
	} 
\end{ex}

% nguồn: Ok_Toan11_TracNghiem_Lan2
\begin{ex}%[1D1N5-2]%[ma:9f04fce4]
	 Có bao nhiêu giá trị nguyên của $m$ để phương trình: $3\sin x+m-1=0$ có nghiệm?
	\choice 
		 {\True $7$}
		{$6$} 
		{$3$} 
		 {$5$}
	\loigiai{ 
		Ta có $3\sin x+m-1=0$ $\Leftrightarrow \sin x=\dfrac{1-m}{3}$, để phương trình có nghiệm ta có \begin{align*}
		&-1\le \dfrac{1-m}{3}  \le 1\\
		 &\Leftrightarrow-2\le m\le 4.
	\end{align*}\\
	Vậy có $7$ giá trị nguyên từ $-2$ đến $4$ thoả mãn  yêu cầu bài toán.
 }
\end{ex}

% nguồn: Ok_Toan11_TracNghiem_Lan2
\begin{ex}%[1D1N5-2]%[ma:d94b5478]
	Có bao nhiêu giá trị  nguyên của tham số $m$ thuộc đoạn $\left[-2018;2018\right]$ để phương trình $m\cos x+1=0$ có nghiệm?
	 \choice
		{\True  $4036$}
		{$4037$} 
		 {$2018$}
		 {$2019$}
	\loigiai{ 
		TH1: Nếu $m=0$ thì phương trình đã cho  vô nghiệm.\\
		TH2: Nếu $m\ne 0$ thì phương trình $m\cos  x+1=0\Leftrightarrow \cos x=-\dfrac{1}{m}$.\\
		Phương trình đã  cho có nghiệm $\Leftrightarrow-1\le-\dfrac{1}{m} \le 1$
		$\Leftrightarrow \hoac{&m\ge 1\\&m\le-1.}$\\ 
		Kết hợp với điều kiện $m$ nguyên và $m$ thuộc đoạn $\left[-2018;2018\right]$ suy ra $m\in  \left\{1;2;3;\ldots;2018\right\}$ hoặc $m\in \left\{-2018;\ldots;-3;-2;-1\right\}$.\\
		Vậy có $4036$  giá trị của tham số $m$ thỏa mãn đề bài.
	 }
\end{ex}

% nguồn: Ok_Toan11_TracNghiem_Lan2
\begin{ex}%[1D1N5-2]%[ma:5f5ff492]
	Tìm tất cả các giá trị của tham số $m$ để  phương trình $\cos ^2x=m-1$ có nghiệm.
	 \choice
		{$m\le  2$}
		{$1< m  < 2$}
		 {$m\ge 1$}
		{\True  $1\le m\le 2$}
	 \loigiai{
		Do  $0\le \cos ^2x\le 1$, $\forall x\in \mathbb{R}$ nên phương trình có nghiệm khi $0\le m-1\le 1\Leftrightarrow 1\le m\le 2$.
	 }
\end{ex}

% nguồn: Ok_Toan11_TracNghiem_Lan2
\begin{ex}%[1D1N5-2]%[ma:14da398d]
	Phương trình  $\sin 5x=m$ có nghiệm khi
	 \choice
		 {$\left|m\right|\le 5$}
		 {$m\le 5$}
		{\True $\left|m\right|\le  1$}
		{$m\le  1$}
	 \loigiai{
		Ta có $-1\le \sin 5x\le 1\Leftrightarrow-1\le m\le 1\Leftrightarrow \left|m\right|\le 1$. 
	 }
\end{ex}

% nguồn: Ok_Toan11_TracNghiem_Lan2
\begin{ex}%[1D1N5-2]%[ma:f791ddb3]
	Tìm tất cả các giá trị thực của tham số  $m$ để phương trình $\sin x=m+1$ có nghiệm?
	\choice 
		{$1\le  m$}
		{$0\le  m\le 1$}
		 {$m\le 0$}
		 {\True $-2\le m\le 0$}
	\loigiai{ 
		Phương trình: $\sin x=m+1$ có  nghiệm $\Leftrightarrow-1\le m+1\le 1\Leftrightarrow-2\le m\le 0$.
	} 
\end{ex}

% nguồn: Ok_Toan11_TracNghiem_Lan2
\begin{ex}%[1D1N5-2]%[ma:6f4e79c4]
	Tìm tất cả các giá trị thực của tham số  $m$ để phương trình $\cos ^2(\dfrac{x}{2}-\dfrac{\pi}{2})=m$ có nghiệm.
	 \choice
		{$-1\le m\le  1$}
		{$m\le 1$} 
		{$m\ge 0$} 
		 {\True $0\le m\le 1$}
	 \loigiai{
		Ta có: $0\le \cos ^2(\dfrac{x}{2}-\dfrac{\pi}{2})\le  1$. Vậy để phương trình có nghiệm thì $0\le m\le 1$.
	 }
\end{ex}

% nguồn: Ok_Toan11_TracNghiem_Lan2
\begin{ex}%[1D1N5-2]%[ma:e938578b]
	Tìm $m$ để phương trình $\sin  3x-6-5m=0$ có nghiệm.
	\choice 
		{\True $-\dfrac{7}{5} \le m\le-1$} 
		{$-\dfrac{7}{5}  < m <-1$}
		{$\hoac{&m\ge-1  \\&m\le-\dfrac{7}{5}}$}
		{$\hoac{&m >-1 \\&m <-\dfrac{7}{5}}$} 
	\loigiai{ 
		Ta có $\sin 3x-6-5m=0\Leftrightarrow  \sin 3x=6+5m$.\\
		Phương trình đã  cho có nghiệm $\Leftrightarrow-1\le 6+5m\le 1 \Leftrightarrow-\dfrac{7}{5} \le m\le-1$.
	} 
\end{ex}

% nguồn: Ok_Toan11_TracNghiem_Lan2
\begin{ex}%[1D1N5-2]%[ma:2239f529]
	Có tất cả bao nhiêu giá trị nguyên của tham số $m$ để phương trình $\sqrt{3} \cos x+m-1=0$ có  nghiệm?
	\choice 
		 {$1$}
		 {$2$}
		{\True $3$} 
		 {Vô số}
	\loigiai{ 
		Ta có $\sqrt{3} \cos x+m-1=0\Leftrightarrow \cos  x=\dfrac{1-m}{\sqrt{3}}$.\\
		 Phương trình có nghiệm $\Leftrightarrow-1\le \dfrac{1-m}{\sqrt{3}} \le 1\Leftrightarrow 1-\sqrt{3} \le m\le 1+\sqrt{3}$.\\
		Do $m \in \mathbb{Z}$ nên  $m\in \left\{0;1;2\right\}$.\\
		Vậy có tất  cả 3 giá trị nguyên của tham số $m$ thoả mãn.
	 }
\end{ex}

% nguồn: Ok_TracNghiemLop11_Lan1
\begin{ex}%[1D1N5-2]%[ma:27f6b57c]
	Phương  trình nào sau đây vô nghiệm?
	\choice 
		 {$\sin x=1$}
		 {$\sin x=-\dfrac{1}{2}$}
		{\True $\sin x=2$} 
		 {$\sin x=0$}
	 \loigiai{
	Phương trình  vô nghiệm là $\sin x=2$}
\end{ex}

% nguồn: Ok_Toan11_TracNghiem_Lan2
\begin{ex}%[1D1H5-2]%[ma:3f918231]
	Trong các phương trình sau, phương trình nào  vô nghiệm?
	 \choice
		{$\tan  x=99$}
		{\True $\cos \left( 2x-\dfrac{\pi}{2}  \right)=\dfrac{2\pi}{3}$}
		 {$\cot 2\,018x=2\,017$}
		{$\sin  2x=-\dfrac{3}{4}$}
	\loigiai{Vì $\dfrac{2\pi}{3}>1$ là nên phương  trình $\cos \left( 2x-\dfrac{\pi}{2} \right)=\dfrac{2\pi}{3}$ vô nghiệm.}
\end{ex}

% nguồn: Ok_TracNghiemLop11_Lan1
\begin{ex}%[1D1H5-2]%[ma:af5d1ada]
	Tìm tất  cả các giá trị thực của tham số $m$ để phương trình $\sin x = m+1$ có nghiệm
	\choice 
		{$-1  \leq m \leq 1$}
		{\True $-2 \leq m  \leq 0$}
		 {$0 \leq m \leq 2$}
		{$0  \leq m \leq 1$}
	\loigiai{ 
		Vì $-1 \leq \sin x \leq 1$ nên phương trình $\sin  x = m+1$ có nghiệm khi và chỉ khi
		$$-1 \leq m+1 \leq 1  \Leftrightarrow -2 \leq m \leq 0.$$
	 }
\end{ex}

% nguồn: Ok_TracNghiemLop11_Lan1
\begin{ex}%[1D1H5-2]%[ma:3674969a]
	 Phương trình $\sin x=m$ vô nghiệm khi và chỉ khi
	 \choice
		 {\True $\hoac{&m<-1\\& m>1}$}
		{$-1\le m\le 1$} 
		 {$m<-1$}
		 {$m>1$}
	\loigiai{Do tập giá trị của hàm số $y=\sin x$ là đoạn $[-1; 1]$, nên phương trình đã  cho vô nghiệm khi và chỉ khi $\hoac{&m<-1\\ &m>1.}$
	} 
\end{ex}

% nguồn: Ok_Toan11_TracNghiem_Lan2
\begin{ex}%[1D1V5-2]%[ma:054d8ca6]
	Số nghiệm của  phương trình $\cos\left(x+\dfrac{\pi}{4}\right)=1$ với $\pi\leq x\leq 5\pi$ là
	\choice 
		{$0$} 
		{$3$} 
		 {$1$}
		 {\True $2$}
	\loigiai{ 
		Ta có  $\cos\left(x+\dfrac{\pi}{4}\right)=1\Leftrightarrow x+\dfrac{\pi}{4}=k2\pi\Leftrightarrow x=-\dfrac{\pi}{4}+k2\pi, k\in\mathbb{Z}$.\\
		Mà  $\pi\leq x\leq 5\pi$ nên $\pi\leq-\dfrac{\pi}{4}+k2\pi\leq 5\pi\Leftrightarrow\dfrac{5}{8}\leq k\leq\dfrac{21}{8}$.\\
		Vì  $k\in\mathbb{Z}$ nên $k\in\left\lbrace1;2\right\rbrace$.\\ Vậy phương trình đã cho có $2$ nghiệm trên $\left[\pi;5\pi\right]$.
	} 
\end{ex}

% nguồn: Ok_Toan11_TracNghiem_Lan2
\begin{ex}%[1D1V5-2]%[ma:485f306e]
	Số nghiệm của phương trình $\sin 2x=\dfrac{\sqrt{3}}{2}$ trong khoảng $\left(0;3\pi\right)$ là 
	 \choice
		 {$4$}
		 {$1$}
		{\True $6$} 
		 {$2$}
	 \loigiai{
		Ta có $\sin 2x=\dfrac{\sqrt{3}}{2}\Leftrightarrow\hoac{&2x=\dfrac{\pi}{3}+k2\pi\\&2x=\dfrac{2\pi}{3}+k2\pi 
		}\Leftrightarrow\hoac{&x=\dfrac{\pi}{6}+k\pi\\ 
		 &x=\dfrac{\pi}{3}+k\pi} (k\in\mathbb{Z})$.\\
		 Mà $x\in\left(0;3\pi\right)$ nên $\hoac{&0<\dfrac{\pi}{6}+k\pi<3\pi\\
		&0<\dfrac{\pi}{3}+k\pi<3\pi 
		}\Leftrightarrow\hoac{&\dfrac{-1}{6}<k<\dfrac{17}{6}\\ 
		&\dfrac{-1}{3}<k<\dfrac{8}{3}.}$\\ 
		 Vì $k\in\mathbb{Z}$ nên $\hoac{&k\in\left\lbrace0;1;2\right\rbrace\\
		 &k\in\left\lbrace0;1;2\right\rbrace.}$\\
		Vậy phương trình $\sin 2x=\dfrac{\sqrt{3}}{2}$ có $6$ nghiệm  trong khoảng $\left(0;3\pi\right)$.
	} 
\end{ex}

% nguồn: Ok_Toan11_TracNghiem_Lan2
\begin{ex}%[1D1V5-2]%[ma:6dcbb524]
	Tổng các nghiệm thuộc khoảng  $\left(0;2\pi\right)$ của phương trình $5\cos x-2=0$ là
	 \choice
		{$S=3\pi$} 
		{\True $S=2\pi$} 
		 {$S=0$}
		{$S=4\pi$} 
	\loigiai{ 
		Ta có $5\cos x-2=0\Leftrightarrow\cos x=\dfrac{2}{5}\Leftrightarrow x=\pm\arccos \left(\dfrac{2}{5}\right)+k2\pi, k\in{\bf\mathbb{Z}}$.\\ 
		Xét trên $\left(0;2\pi\right)$ phương trình có hai  nghiệm $x=\arccos\left(\dfrac{2}{5}\right)$ và $x=-\arccos\left(\dfrac{2}{5}\right)+2\pi$.\\
		Do vậy tổng tất cả các nghiệm  của phương trình bằng $\arccos\left(\dfrac{2}{5}\right)-\arccos\left(\dfrac{2}{5}\right)+2\pi=2\pi$.
	} 
\end{ex}

% nguồn: Ok_Toan11_TracNghiem_Lan2
\begin{ex}%[1D1V5-2]%[ma:f416d87f]
	Phương trình $\sin 2x=\dfrac{-1}{2}$ có bao  nhiêu nghiệm thoả mãn $0<x<\pi$?
	\choice 
		{$1$} 
		 {$3$}
		{\True $2$} 
		 {$4$}
	\loigiai{ 
		Ta có $\sin  2x=\dfrac{-1}{2}\Leftrightarrow\hoac{&2x=\dfrac{-\pi}{6}+k2\pi\\
		&2x=\dfrac{7\pi}{6}+k2\pi}\left(k\in{\bf\mathbb{Z}}\right)\Leftrightarrow\hoac{&x=\dfrac{-\pi}{12}+k\pi\\ 
		 &x=\dfrac{7\pi}{12}+k\pi}\left(k\in{\bf\mathbb{Z}}\right)$.\\
		\begin{itemize}
			 \item \textbf{Trường hợp 1:} $x=\dfrac{-\pi}{12}+k\pi\left(k\in{\bf\mathbb{Z}}\right)$.\\
			Do $0<x<\pi\Rightarrow 0<\dfrac{-\pi}{12}+k\pi<\pi\Rightarrow\dfrac{1}{12}<k<\dfrac{13}{12}$.\\ 
			Vì  $k\in\mathbb{Z}$ nên $k=1$ suy ra $x=\dfrac{11\pi}{12}$.
			\item \textbf{Trường hợp 2:}  $x=\dfrac{7\pi}{12}+k\pi\left(k\in{\bf\mathbb{Z}}\right)$.\\
			Do  $0<x<\pi\Rightarrow 0<\dfrac{7\pi}{12}+k\pi<\pi\Rightarrow\dfrac{-7}{12}<k<\dfrac{5}{12}$.\\
			Vì $k\in\mathbb{Z}$ nên $k=0$ suy ra  $x=\dfrac{7\pi}{12}$.
		\end{itemize}
		Vậy phương trình $\sin 2x=\dfrac{-1}{2}$ có $2$  nghiệm thoả mãn $0<x<\pi$.
	 }
\end{ex}

% nguồn: Ok_Toan11_TracNghiem_Lan2
\begin{ex}%[1D1V5-2]%[ma:93682c90]
	Số nghiệm của phương trình  $\sin\left(5x-\dfrac{\pi}{3}\right)=1$ trên $\left[0;2\pi\right]$ là
	 \choice
		 {\True $5$}
		{$3$} 
		 {$4$}
		{$2$} 
	\loigiai{ 
		Ta có $\sin\left(5x-\dfrac{\pi}{3}\right)=1\Leftrightarrow  5x-\dfrac{\pi}{3}=\dfrac{\pi}{2}+k2\pi\Leftrightarrow x=\dfrac{\pi}{6}+\dfrac{k2\pi}{5}, k\in\mathbb{Z}$.\\
		Khi đó $x\in\left[0;2\pi\right]\Leftrightarrow 0\leq\dfrac{\pi}{6}+\dfrac{k2\pi}{5}\leq 2\pi\Leftrightarrow\dfrac{-5}{12}\leq k\leq\dfrac{55}{12}$.\\ 
		Vì  $k\in\mathbb{Z}$ nên $k\in \left\lbrace 0;1;2;3;4\right\rbrace$.\\
		Vậy có $5$ nghiệm của phương trình $\sin\left(5x-\dfrac{\pi}{3}\right)=1$ thuộc khoảng $\left[0;2\pi\right]$. 
	} 
\end{ex}

% nguồn: Ok_Toan11_TracNghiem_Lan2
\begin{ex}%[1D1V5-2]%[ma:8cf30ba2]
	Tổng các nghiệm dương nhỏ nhất và nghiệm âm lớn nhất của phương  trình $\cos\left(2x-\dfrac{\pi}{6}\right)=\dfrac{\sqrt{3}}{2}$ là
	 \choice
		 {$\dfrac{\pi}{12}$}
		{\True  $\dfrac{-2\pi}{3}$}
		{$\dfrac{\pi}{4}$} 
		 {$\dfrac{-3\pi}{4}$}
	\loigiai{ 
		Ta  có
		\[\cos\left(2x-\dfrac{\pi}{6}\right)=\dfrac{\sqrt{3}}{2} 
		\Leftrightarrow  \cos\left(2x-\dfrac{\pi}{6}\right)=\cos\left(\dfrac{\pi}{6}\right)
		\Leftrightarrow \hoac{&2x-\dfrac{\pi}{6}=\dfrac{\pi}{6}+k2\pi \\&2x-\dfrac{\pi}{6}=-\dfrac{\pi}{6}+k2\pi} 
		 \Leftrightarrow \hoac{&x=\dfrac{\pi}{6}+k\pi \\&x=k\pi} (k\in\mathbb{Z}).\]
		\textbf{Tìm nghiệm dương  nhỏ nhất:}
		\begin{itemize}
			\item Với họ nghiệm $x=k\pi$: để $x>0$ thì $k>0$.\\ Giá trị $k$ nguyên  nhỏ nhất là $k=1$, ta được $x=\pi$.
			\item Với họ nghiệm $x=\dfrac{\pi}{6}+k\pi$: để $x>0$ thì $k\pi > -\dfrac{\pi}{6} \Leftrightarrow k  > -\dfrac{1}{6}$. \\
			Giá trị $k$ nguyên nhỏ nhất  là $k=0$, ta được $x=\dfrac{\pi}{6}$.
		\end{itemize}
		So sánh hai giá trị $\pi$ và $\dfrac{\pi}{6}$, ta thấy nghiệm dương nhỏ nhất của  phương trình là $x=\dfrac{\pi}{6}$.\\
		\textbf{Tìm nghiệm âm lớn nhất:} 
		\begin{itemize}
			\item Với họ nghiệm $x=k\pi$: để $x<0$ thì $k<0$.\\ Giá trị $k$  nguyên lớn nhất là $k=-1$, ta được $x=-\pi$.
			 \item Với họ nghiệm $x=\dfrac{\pi}{6}+k\pi$: để $x<0$ thì $k\pi < -\dfrac{\pi}{6} \Leftrightarrow k < -\dfrac{1}{6}$.\\
			Giá trị  $k$ nguyên lớn nhất là $k=-1$, ta được $x=\dfrac{\pi}{6}-\pi = -\dfrac{5\pi}{6}$.
		\end{itemize}
		So sánh hai giá trị $-\pi$  và $-\dfrac{5\pi}{6}$, ta thấy nghiệm âm lớn nhất của phương trình là $x=-\dfrac{5\pi}{6}$.
		\\Tổng của  nghiệm dương nhỏ nhất và nghiệm âm lớn nhất là
		$S  = \dfrac{\pi}{6} -\dfrac{5\pi}{6}= -\dfrac{2\pi}{3}$.
	 }
\end{ex}

% nguồn: Ok_Toan11_TracNghiem_Lan2
\begin{ex}%[1D1V5-2]%[ma:e3c32d21]
	Tổng các nghiệm của phương trình $\tan 2x=\tan x$ trên $\left[-\pi;2\pi\right]$ bằng 
	 \choice
		 {$\pi$}
		{$\dfrac{\pi}{2}$} 
		 {$4\pi$}
		{\True $2\pi$} 
	\loigiai{ 
		Điều  kiện xác định
		$\heva{&\cos 2x\neq  0\\&\cos x\neq 0}
		\Leftrightarrow\heva{&x\neq\dfrac{\pi}{4}+\dfrac{k\pi}{2}\\ 
		&x\neq\dfrac{\pi}{2}+k\pi} (k\in\mathbb{Z})$.\\ 
		Khi đó $\tan 2x=\tan x\Leftrightarrow 2x=x+k\pi\Leftrightarrow x=k\pi, k\in\mathbb{Z}$.\\ Do $x\in\left[-\pi;2\pi\right]$  nên $x\in\left\lbrace-\pi;0;\pi;2\pi\right\rbrace$.\\
		Vậy tổng các nghiệm của phương trình $\tan 2x=\tan x$ trên $\left[-\pi;2\pi\right]$ bằng  $2\pi$.
	 }
\end{ex}


%%%%%%%%%%%%%%%% ĐÚNG/SAI %%%%%%%%%%%%%%%%

% nguồn: Ok_dungsai_lop11_Lan1
\begin{ex}%[1D1N5-2]%[ma:a0ad4965]
 Phương  trình nào dưới đây có nghiệm.
  \choiceTF
  {\True $\sin  x=\dfrac{2017}{2018}$}
  {\True  $\tan x=\dfrac{2018}{2017}$}
   {$\cos x=\dfrac{2018}{2017}$}
   {\True $\cot x=\dfrac{2017}{2018}$}
 \loigiai{ 
  Do  $\dfrac{2018}{2017}>1$ nên $\cos x=\dfrac{2018}{2017}$ vô nghiệm, các phương trình còn lại đều có nghiệm.
 } 
\end{ex}

% nguồn: Ok_dungsai_lop11_Lan1
\begin{ex}%[1D1H5-2]%[ma:91801580]
  [1D1H5-5]
 Phương  trình $\sin 3 x=\sin x$. Xét tính đúng - sai của các phát biểu sau.
  \choiceTF
   {Tập nghiệm của phương trình là $S=\left\{\dfrac{\pi}{4}+k \pi \bigg| k \in \mathbb{Z}\right\}$}
  {\True  $x=0$ là một nghiệm của phương trình}
  {\True Nghiệm dương nhỏ nhất  của phương trình là $x=\dfrac{\pi}{4}$}
  {\True Tổng các nghiệm thuộc đoạn  $\left[-\pi; \dfrac{5 \pi}{2}\right]$ của phương trình bằng $\dfrac{29 \pi}{4}$}
 \loigiai{ 
  \begin{itemchoice}
   \itemch Sai. $\sin 3 x=\sin x \Leftrightarrow\hoac{&3 x=x+k 2 \pi \\ &3 x=\pi-x+k 2 \pi } \Leftrightarrow\hoac{&x=k \pi \\ &x=\dfrac{\pi}{4}+\dfrac{k \pi}{2}},  (k \in \mathbb{Z})$.
   \itemch Đúng. Với  $x=0$ ta có $\sin (3\cdot 0)=\sin 0$.
   \itemch Đúng. Với $x=k \pi$, $(k \in \mathbb{Z})$ thì nghiệm dương nhỏ nhất là  $\pi$.\\
   Với $x=\dfrac{\pi}{4}+\dfrac{k \pi}{2}(k \in \mathbb{Z})$ thì nghiệm dương nhỏ nhất là $\dfrac{\pi}{4}$. 
   \itemch Với $x=k \pi$, $(k \in \mathbb{Z})$, các nghiệm thuộc đoạn $\left[-\pi; \dfrac{5 \pi}{2}\right]$  là $x \in\{-\pi; 0; \pi; 2 \pi\}$.\\
   Với $ x=\dfrac{\pi}{4}+\dfrac{k \pi}{2}$, $(k \in \mathbb{Z})$, các nghiệm thuộc đoạn  $\left[-\pi; \dfrac{5 \pi}{2}\right]$ là $x \in\left\{\dfrac{-3 \pi}{4}; \dfrac{-\pi}{4}; \dfrac{\pi}{4}; \dfrac{3 \pi}{4}; \dfrac{5 \pi}{4}; \dfrac{7 \pi}{4}; \dfrac{9 \pi}{4}\right\}$.\\
   Tổng các nghiệm thuộc đoạn $\left[-\pi; \dfrac{5 \pi}{2}\right]$ của phương trình  là $$-\pi+0+\pi+2 \pi+\dfrac{-3 \pi}{4}+\dfrac{-\pi}{4}+\dfrac{\pi}{4}+\dfrac{3 \pi}{4}+\dfrac{5 \pi}{4}+\dfrac{7 \pi}{4}+\dfrac{9 \pi}{4}=\dfrac{29 \pi}{4}.$$
  \end{itemchoice}
  }
\end{ex}

% nguồn: Ok_dungsai_lop11_Lan1
\begin{ex}%[1D1H5-2]%[ma:d98dabca]
  Xét tính đúng sai của các mệnh đề sau.
  \choiceTF
  {Phương trình $\sin  x=\sin m$ có nghiệm khi $m\in [-1;1]$}
  {Phương trình $\cos 2x=\cos  m$ có nghiệm khi $m\in\left[-\dfrac{1}{2};\dfrac{1}{2}\right]$}
  {\True Phương trình  $\cos x=\cos m$ có nghiệm với mọi $m\in\mathbb{R}$}
  {\True Phương trình $\sin 2x=\sin  m$ có nghiệm với mọi $m\in\mathbb{R}$}
 \loigiai{Do $\sin m$ và $\cos m$ đều là các số thuộc đoạn $[-1;1]$ nên  phương trình $\cos x=\cos m$ và $\sin 2x=\sin m$ đều có nghiệm với mọi $m\in\mathbb{R}$.}
\end{ex}

% nguồn: Ok_dungsai_lop11_Lan1
\begin{ex}%[1D1H5-2]%[ma:1fd4f093]
 Cho phương trình $\sin  x=m$. Xét tính đúng - sai của các mệnh đề sau.
 \choiceTF 
  { Với $m=0$ thì phương trình có tập nghiệm là $x=k 2 \pi$, $(k \in \mathbb{Z})$} 
  {\True Với $-1 \leq m  \leq 1$ thì phương trình có nghiệm}
  {\True Với $m=2024$  thì phương trình vô nghiệm}
  { Với $m=1$ thì phương  trình có tập nghiệm là $x=\dfrac{\pi}{2}+k 360^{\circ}$, $(k \in \mathbb{Z})$}
  \loigiai{
  \begin{itemchoice}
   \itemch Sai. Vì $\sin x=0  \Leftrightarrow x=k \pi$, $(k \in \mathbb{Z})$.
   \itemch Đúng. Vì đường thẳng $y=m$, $(-1  \leq m \leq 1)$ luôn cắt đồ thị hàm số $y=\sin x$ nên phương trình đã cho có nghiệm.
   \itemch Đúng. Vì $2024 \notin[-1; 1]$ nên  thì phương trình đã cho vô nghiệm.
   \itemch Sai. Vì $\sin x=1  \Leftrightarrow x=\dfrac{\pi}{2}+k 2 \pi$, $(k \in \mathbb{Z})$.
  \end{itemchoice}
 } 
\end{ex}

% nguồn: Ok_dungsai_lop11_Lan1
\begin{ex}%[1D1H5-2]%[ma:46bccfcd]
 Xét tính đúng sai của  các mệnh đề sau.
  \choiceTF
  {\True Phương trình  $\cos x=\cos m$ có nghiệm với mọi $m\in\mathbb{R}$}
  {Phương  trình $\sin x=\sin m$ có nghiệm khi $m\in [-1;1]$}
  {Phương trình $\cos 2x=\cos  m$ có nghiệm khi $m\in\left[-\dfrac{1}{2};\dfrac{1}{2}\right]$}
  {\True Phương trình $\sin 2x=\sin  m$ có nghiệm với mọi $m\in\mathbb{R}$}
 \loigiai{Do  $\sin m$ và $\cos m$ đều là các số thuộc đoạn $[-1;1]$ nên phương trình $\cos x=\cos m$ và $\sin 2x=\sin m$ đều có nghiệm với mọi $m\in\mathbb{R}$.}
\end{ex}

% nguồn: Ok_dungsai_lop11_Lan1
\begin{ex}%[1D1V5-2]%[ma:00acbcc2]
 Cho hàm  số $y=\sin x$. Xét tính đúng-sai của các phát biểu sau:
  \choiceTF
  {Tập xác định của hàm số  $\mathbb{R}$ và tập giá trị của hàm số là $T=\left[0;1\right]$}
  {\True  Đồ thị hàm số đi qua điểm $\left(\dfrac{\pi}{4}; \dfrac{\sqrt{2}}{2} \right)$}
  {\True Đồ  thị hàm số $y=\sin x$ cắt trục hoành tại vô số điểm}
  {Đồ thị hàm số $y=\sin x$ không cắt đường  thẳng $y=2m+3$ (với $m\in \mathbb{R}$ là tham số) khi $m < \dfrac{-3}{2}$}
  \loigiai{
  \begin{itemchoice}
    \itemch Hàm số $y=\sin x$ có\\
   Tập  xác định là $\mathscr{D}=\mathbb{R}$.\\
   Tập giá trị là $T=\left[-1;1\right]$. 
   \itemch Với  $x=\dfrac{\pi}{4}$.\\
   Ta có $y=\sin x=\sin \dfrac{\pi}{4}=\dfrac{\sqrt{2}}{2}$. 
   \itemch Đồ thị  hàm số $y=\sin x$ cắt trục hoành tại vô số điểm.\\
   Hàm số $y=\sin x$ và trục hoành $y=0$.\\ 
   Phương trình hoành  độ số giao điểm của hai đồ thị là: $\sin x=0\Leftrightarrow x=k\pi \left(k\in \mathbb{Z}\right)$.\\
   Với mỗi giá trị  $k\in \mathbb{R}$ ta được một nghiệm tức là một giao điểm. Phương trình trên có vô số nghiệm.\\
   Vậy đồ thị hàm số $y=\sin x$ cắt  trục hoành tại vô số điểm.
   \itemch Phương trình hoành độ giao điểm đồ thị hàm số $y=\sin x$ và đường thẳng $y=2m+3$ là $\sin  x=2m+3$.\\
   Đồ  thị hàm số $y=\sin x$ không cắt đường thẳng $y=2m+3$ (với $m\in \mathbb{R}$ là tham số) khi phương trình $\sin x=2m+3$ vô nghiệm.\\
   Vậy ta có  $\hoac{&2m+3> 1\\&2m+3<-1}\Leftrightarrow \hoac{&m >-1\\&m <-2}\Leftrightarrow m\in \left(-\infty;-2\right)\cup \left(-1;-\infty \right)$.
  \end{itemchoice}
  }
\end{ex}

% nguồn: Ok_dungsai_lop11_Lan1
\begin{ex}%[1D1V5-2]%[ma:230065f8]
 Cho phương trình $\sin 2 x=m$. Xét tính đúng - sai của các phát  biểu sau.
  \choiceTF
  {\True Với $m=0$, phương trình có nghiệm là $x=\dfrac{k \pi}{2}$} 
  {Phương trình có nghiệm khi $m \in(-1; 1)$} 
  { Phương trình  có nghiệm là $x=\dfrac{\pi}{3}$ khi $m=\dfrac{-1}{2}$}
  {\True Phương trình có hai nghiệm phân biệt trên $\left(0; \dfrac{\pi}{2}\right)$ khi  $m \in(0; 1)$}
  \loigiai{
  \begin{itemchoice}
   \itemch Đúng. Với $m=0$, ta có  $\sin 2 x=0 \Leftrightarrow 2 x=k \pi \Leftrightarrow x=\dfrac{k \pi}{2}$.
    \itemch Phương trình có nghiệm khi $m \in[-1; 1]$.
   \itemch Sai. Phương trình có nghiệm là $x=\dfrac{\pi}{3}$ khi $m=\sin \dfrac{2 \pi}{3}=\dfrac{\sqrt{3}}{2}$. 
   \itemch Đúng. Ta có $x \in\left(0; \dfrac{\pi}{2}\right)$ nên $2 x \in(0; \pi)$. Dưạ vào đường tròn lượng giác ta thấy phương trình có $\sin 2 x=m$ hai nghiệm phân biệt trên $\left(0; \dfrac{\pi}{2}\right)$ khi $m \in(0; 1)$. 
  \end{itemchoice}
  }
\end{ex}


%%%%%%%%%%%%%%%% TRẢ LỜI NGẮN %%%%%%%%%%%%%%%%

% nguồn: Ok_TraLoiNgan_Lop11_Lan1
\begin{ex}%[1D1H5-2]%[ma:4b3e787e]
	Có bao  nhiêu giá trị nguyên của $m\in[0;2024]$ để hàm số $y=\sin \left(\dfrac{2x}{x^2+2x+m^2+3} \right)$ có tập xác định là $\mathscr{D}=\mathbb{R}$?
	\shortans{2025} 
	\loigiai{ 
		Điều kiện $x^2+2x+m^2+3\ne 0$.\\ 
		Ta có  $x^2+2x+m^2+3=(x^2+2x+1)+m^2+2=(x+1)^2+m^2+2\ge m^2+2>0, \forall m\in \mathbb{R}$.
		Vậy có  $2025$ giá trị nguyên của $m$ thỏa mãn.
	 }
\end{ex}

% nguồn: Ok_TraLoiNgan_Lop11_Lan2
\begin{ex}%[1D1H5-2]%[ma:6484403a]
	Gọi $S$ là tập hợp tất cả các giá trị nguyên của  tham số $m$ để phương trình \break $\cos\left(2x-\dfrac{\pi}{3}\right)-m=2$ có nghiệm. Tổng của các phần tử trong $S$ bằng bao nhiêu?
	 \shortans{-6}
	\loigiai{ 
		 Phương trình $\cos\left(2x-\dfrac{\pi}{3}\right)-m=2\Leftrightarrow\cos\left(2x-\dfrac{\pi}{3}\right)=m+2$.\\
		Phương trình có nghiệm $\Leftrightarrow-1\leq  m+2\leq 1\Leftrightarrow-3\leq m\leq-1$.\\
	 Mà $m\in\mathbb{Z}$ nên $S=\{-3;-2;-1\}$. Do đó $T=(-3)+(-2)+(-1)=-6$.}
\end{ex}

% nguồn: Ok_TraLoiNgan_Lop11_Lan2
\begin{ex}%[1D1H5-2]%[ma:d36dd699]
	Tập hợp các giá trị của tham số $m$ để phương trình $\sin 2x+2=m$ có nghiệm là $\left[a;b\right]$. Khi đó $a+b$ bằng bao  nhiêu?
	\shortans{4} 
	 \loigiai{
		Ta có $\sin2x+2=m\Leftrightarrow\sin2x=m-2$  có nghiệm khi và chỉ khi $-1\le m-2\le 1\Leftrightarrow 1\le m\le 3\Leftrightarrow m\in\left[1;3\right]$.\\
		Vậy $a+b=4$ 
	} 
\end{ex}

% nguồn: Ok_TraLoiNgan_Lop11_Lan2
\begin{ex}%[1D1H5-2]%[ma:8058acd0]
	Tập hợp các giá trị của tham số $m$ để  phương trình $3\sin^2x=2m-1$ có nghiệm là $\left[a;b\right]$. Khi đó $2a+b$ bằng bao nhiêu?
	\shortans{3} 
	 \loigiai{
		Ta  có $3\sin^2x=2m-1\Leftrightarrow{\sin^2}x=\dfrac{2m-1}{3}$\\
		Để phương trình có nghiệm  thì $\heva{ &\dfrac{2m-1}{3}\le 1\\ &\dfrac{2m-1}{3}\ge 0}\Leftrightarrow\heva{& m\le 2\\& m\ge\dfrac{1}{2}}\Leftrightarrow\dfrac{1}{2}\le m\le 2$.\\
		Từ đó có  $a=\dfrac{1}{2}$ , $b=2$.\\
		Vậy $2a+b=2\cdot \dfrac{1}{2}+2=3$. 
	 }
\end{ex}


%%%%%%%%%%%%%%%% TỰ LUẬN %%%%%%%%%%%%%%%%

% nguồn: Ok_TuanLuan_Toan11_Lan2
\begin{ex}%[1D1H5-2]%[ma:2a4980f6]
	 Giải các phương trình lượng giác sau
	\begin{multicols}{2}
		\begin{enumerate}
			 \item $\sin\left(2x+\dfrac{2\pi}{5}\right)=0$ với $x\in\left(\dfrac{\pi}{2};\dfrac{3\pi}{2}\right)$;
			\item  $\tan\left(\dfrac{x}{2}+\dfrac{\pi}{6}\right)=-1$ với $x\in\left(-\dfrac{\pi}{2};\dfrac{\pi}{2}\right)$.
		\end{enumerate}
	\end{multicols}
	\loigiai{ 
		\begin{enumerate}
			\item Ta  có $\sin\left(2x+\dfrac{2\pi}{5}\right)=0\Leftrightarrow 2x+\dfrac{2\pi}{5}=k\pi\Leftrightarrow x=-\dfrac{\pi}{5}+\dfrac{k\pi}{2}$.
			\\ 
			 Do $x\in\left(\dfrac{\pi}{2};\dfrac{3\pi}{2}\right)$ nên $\dfrac{\pi}{2}<-\dfrac{\pi}{5}+\dfrac{k\pi}{2}<\dfrac{3\pi}{2}\Leftrightarrow\dfrac{7}{5}<k<\dfrac{17}{5}$.\\
			 Vì $k\in\mathbb{Z}$ nên $k\in\left\lbrace 2;3\right\rbrace $.\\
			Vậy $x\in\left\lbrace\dfrac{4\pi}{5};\dfrac{13\pi}{10}\right\rbrace$. 
			\item Ta có $\tan\left(\dfrac{x}{2}+\dfrac{\pi}{6}\right)=-1\Leftrightarrow\dfrac{x}{2}+\dfrac{\pi}{6}=-\dfrac{\pi}{4}+k\pi\Leftrightarrow x=-\dfrac{5\pi}{6}+k2\pi$.\\ 
			Do  $x\in\left(-\dfrac{\pi}{2};\dfrac{\pi}{2}\right)$ nên $-\dfrac{\pi}{2}<-\dfrac{5\pi}{6}+k2\pi<\dfrac{\pi}{2}\Leftrightarrow\dfrac{1}{6}<k<\dfrac{2}{3}$.\\
			Mặt khác $k\in\mathbb{Z}$ nên không có giá trị nào của  $k$ thoả mãn.\\
			Vậy  phương trình không có nghiệm trên khoảng $\left(-\dfrac{\pi}{2};\dfrac{\pi}{2}\right)$.
		\end{enumerate}
	 }
\end{ex}

% nguồn: Ok_TuanLuan_Toan11_Lan2
\begin{ex}%[1D1H5-2]%[ma:ee8717a7]
	Tìm $m$  để phương trình $\sqrt{2}\sin\left(x+\dfrac{\pi}{4}\right)=m$ có nghiệm $x\in\left(0;\dfrac{\pi}{2}\right)$.
	 \loigiai{
		 Ta có $0<x<\dfrac{\pi}{2}\Rightarrow\dfrac{\pi}{4}<x+\dfrac{\pi}{4}<\dfrac{3\pi}{4}$ $\Rightarrow\dfrac{\sqrt{2}}{2}<\sin\left(x+\dfrac{\pi}{4}\right)\le 1$.\\
		 Phương trình đã cho có nghiệm $x\in\left(0;\dfrac{\pi}{2}\right)$ khi $\dfrac{\sqrt{2}}{2}<\dfrac{m}{\sqrt{2}}\le 1\Leftrightarrow 1<m\le\sqrt{2}$.
	} 
\end{ex}
